\begin{trapbox}

	\begin{description}[style=nextline,leftmargin=2.8em,labelsep=0.5em,itemsep=0.35em]
		\item[\textbf{\textcolor{CTrap}{易错点一（忽略λ范围）：}}]
			学生使用曲线系时不检查 $\lambda<a^{2}$ 且 $\lambda\neq b^{2}$，导致方程无意义。
		\item[\textbf{\textcolor{CTrap}{正确理解（先验条件）：}}]
			必须首先确认 $\lambda$ 满足条件，否则曲线系不成立。
		\item[\textbf{\textcolor{CTrap}{错因分析（思维盲区）：}}]
			习惯认为方程总是椭圆，忽略参数限制。
	\end{description}

	\medskip
	\begin{description}[style=nextline,leftmargin=2.8em,labelsep=0.5em,itemsep=0.35em]
		\item[\textbf{\textcolor{CTrap}{易错点二（焦点偏移错觉）：}}]
			误以为焦点会随 $\lambda$ 改变而移动。
		\item[\textbf{\textcolor{CTrap}{正确理解（焦点恒定）：}}]
			所有曲线焦点均为 $(\pm c,0)$，焦距恒为 $c=\sqrt{a^{2}-b^{2}}$。
		\item[\textbf{\textcolor{CTrap}{错因分析（类比不当）：}}]
			误认为分母变化导致曲线形状改变的同时焦点也变。
	\end{description}

	\medskip
	\begin{description}[style=nextline,leftmargin=2.8em,labelsep=0.5em,itemsep=0.35em]
		\item[\textbf{\textcolor{CTrap}{易错点三（类型判定混淆）：}}]
			只根据分母是否都正判定椭圆，忽略双曲线可能。
		\item[\textbf{\textcolor{CTrap}{正确理解（分类条件）：}}]
			$\lambda<b^{2}$ 为椭圆，$b^{2}<\lambda<a^{2}$ 为双曲线。
		\item[\textbf{\textcolor{CTrap}{错因分析（条件遗漏）：}}]
			未考虑 $b^{2}<\lambda<a^{2}$ 时 $b^{2}-\lambda<0$，曲线类型为双曲线。
	\end{description}

\end{trapbox}
